\begin{figure}
    \begin{lstlisting}[]
void spmm(row_ptrs: mref<? x idx>, // adj matrix CSR pointers
          col_idxs: mref<? x idx>, // adj matrix CSR indices
          nnz_vals: mref<? x f32>, // matrix nonzero values
          emb_vals: mref<? x ? x f32>, // embedding table
          out_vals: mref<? x ? x f32>){ // output tensor
  // Go through all rows
  for (idx mtx_row_idx=0; mtx_row_idx<n_rows; mtx_row_idx++){
    idx mtx_ptr_beg = row_ptrs[mtx_row_idx];
    idx mtx_ptr_end = row_ptrs[mtx_row_idx+1];
    // Go through all nonzeros in a row
    for (idx mtx_ptr=mtx_ptr_beg; mtx_ptr<mtx_ptr_end; mtx_ptr++){
      f32 mtx_nnz_val = mtx_nnz_vals[mtx_ptr];
      idx mtx_col_idx = mtx_col_idxs[mtx_ptr];
      // Go through all elements in embedding vector
      for (idx emb_col_idx=0; emb_col_idx<emb_len; emb_col_idx++){
        f32 emb_val = emb_vals[mtx_col_idx,emb_col_idx];
        out_vals[mtx_row_idx,emb_col_idx] += mtx_nnz_val * emb_val; }}}}
    \end{lstlisting}
    \vspace{-1em}
    \caption[Pseudocode of graph convolution.]{Pseudocode for the SpMM-style graph convolution in \Cref{fig:gnn-lookup}. The outer loop iterates over graph nodes, the middle loop scans the nonzeros that identify each node's neighbors, and the inner loop traverses the embedding vector fetched for each neighbor. Compared with SpMV, the additional inner loop introduces more control-flow operations and a more regular structure that is generally leveraged for vectorization.} %This operation is similar to CSR sparse-dense matrix multiplication~\cite{aktulga2014ipdps}, although embedding vectors are not rescaled by the matrix values.}
    \label{fig:gc-code}
\end{figure}
